\section{Why the tree of line-bundle orbits does not deform}
\label{sec:tree-extension-obstruction}

This section shows that the tree used to obtain simplicity in
Theorem~\ref{thm:simple-hodge-bundle} obstructs deformation in at
least one polarized Weil direction.  The obstruction concerns the
underlying bundle, even when its line subbundles are allowed to
disappear under deformation.

Retain the full permutation orbits of the simple Hermitian spectrum
$(1,-1,\ldots,n,-n)$ used there.  Within each orbit retain all
transposition edges.  Between distinct orbits retain exactly one
rank-two edge for each edge of a tree.  One may also tensor the
line bundles just constructed by arbitrary elements of
$\operatorname{Pic}^0(A)$.

\begin{theorem}\label{thm:tree-extension-obstruction}
Let $n\ge2$.  Form an ordinary extension
\eqref{eq:positive-bundle-extension} using positive multiplicities
at all vertices and arbitrary extension matrices supported on this
graph.  If it deforms to first order in every polarized Weil tangent
direction, its multiplicities are constant on each full permutation
orbit.  In particular its Chern character is a polynomial in $\theta$.
Consequently an extension with character
$R F_n(N\theta)+cw$, $c\ne0$, cannot be strongly semiregular.
\end{theorem}

The graph hypothesis is essential to the conclusion proved here.
Several bridges between the same two orbit sets, cycles between
orbit sets, or additional nonzero extension entries are not covered.
The proof gives a necessary trace equation for deformation and then
solves that equation on a leaf of the tree.

\subsection{The trace equation allows unfiltered deformations}

Split the underlying smooth bundle as the sum of its line factors.
Its Dolbeault operator is $\bar\partial_0+e$, where $e$ has entries
only from $P$ to $Q$, so $e^2=0$ and $\bar\partial_0e=0$.
Choose a connection whose $(1,0)$ part is block diagonal on these
factors.  Its $(1,1)$ curvature has the line curvatures on the
diagonal and other entries only from $P$ to $Q$.

For a complex-structure deformation $\kappa$, a first-order change
$x$ in the Dolbeault operator must satisfy
\begin{equation}
 \bar\partial_0x+[e,x]+\kappa\lrcorner F^{1,1}=0,
 \label{eq:tree-atiyah}
\end{equation}
up to the common sign convention.  This is obtained by expanding
the square of the deformed operator; its cohomology class is the
Atiyah obstruction.  In the reverse block from $Q$ to $P$ both the
curvature term and $[e,x]$ vanish.  Thus these reverse entries of
$x$ are closed and give arbitrary classes in the corresponding
$\operatorname{Ext}^1$ groups.

Take the trace of the diagonal equation at vertex $v$, and pass
to $H^2(\mathcal O_A)$.  Its prescribed term is
$m_v\kappa\lrcorner c_1(L_v)$.
An edge contribution is the cup product of its extension entry
with a reverse class.  By \eqref{eq:edge-surface-line} it lies in
the one-dimensional line $q_e$ from the associated surface quotient.
This remains true with arbitrary flat twists.  Indeed an isogeny
from a product of an abelian surface and the kernel abelian variety
pulls the difference bundle back to $M\boxtimes\alpha$, where
$M$ has index one and $\alpha\in\operatorname{Pic}^0$.
If $\alpha$ is nontrivial its cohomology vanishes: a unitary
character description and Fourier expansion of the Dolbeault
complex give nonzero Koszul differentials at every frequency.
If $\alpha$ is trivial, K\"unneth and $H^0(M)=H^0(M^{-1})=0$
put both degree-one groups on the surface factor.  Their cup
product therefore lies in its $H^2(\mathcal O)$ line.
Isogeny pullback is injective on line-bundle cohomology by trace,
and is an isomorphism on $H^2(\mathcal O)$, proving the claim.
The two endpoint traces are opposite: interchanging the two
degree-one forms changes the sign.
Therefore a necessary condition for deformation is the existence
of scalar edge flows with
\begin{equation}
 \bigl(m_v\kappa\lrcorner c_1(L_v)\bigr)_v
  =\sum_e f_e\,(q_e\text{ at one endpoint},
                         -q_e\text{ at the other}).
 \label{eq:tree-incidence}
\end{equation}
Allowing arbitrary flows can only enlarge the possible right side.
The diagonal entries of $x$ contribute exact terms and disappear
in cohomology.  Reverse blocks were included throughout, so this
condition imposes no requirement that the deformation preserve
the original filtration.  There is also no omitted higher-product
term: the Dolbeault equation is the full first-order equation.

\subsection{Two facts about the Weil tangent space}

Using the product polarization, the polarized Weil tangent space
is congruent to
\begin{equation}
 \mathcal T=
 \left\{\begin{pmatrix}0&C\\C^{\mathsf t}&0\end{pmatrix}:
                    C\in\operatorname{Mat}_n(\mathbf C)\right\}
 \subset\operatorname{Sym}^2(\mathbf C^{2n})^*.
 \label{eq:weil-offblock-tangent}
\end{equation}
It has dimension $n^2$.  The following arguments work in every
eigenbasis, without a genericity assumption.

First, join coordinate indices $j,k$ when
$\kappa\mapsto\kappa_{jk}$ is not identically zero on $\mathcal T$.
This coordinate graph is connected.  Otherwise the coordinate
partition gives a nontrivial decomposition $W=A_1\oplus A_2$
orthogonal for every form in \eqref{eq:weil-offblock-tangent}.
Write $W=U\oplus V$ in that equation.  For $x=u+v$ and
$y=u'+v'$, simultaneous orthogonality is equivalent to
\[
 u\otimes v'+u'\otimes v=0.
\]
If $u,v$ are both nonzero, its solutions in $y$ form the line
$\mathbf C(u-v)$.  If $A_1$ contains such an $x$, then $A_2$ is
a mixed line, and applying the same argument in reverse makes
$A_1$ a line.  This contradicts $2n\ge4$.
If $A_1$ contains no mixed vector, its linearity puts it in either
$U$ or $V$.  A nonzero vector in $U$, for example, forces
$A_2\subset U$, again contradicting the direct sum.  The graph
is connected as claimed.

Second, the off-diagonal projection of $\mathcal T$ has rank at
least $n^2-n\ge2$.  Every matrix in
\eqref{eq:weil-offblock-tangent} has even rank, and congruence
preserves rank.  Let $S$ be the intersection with diagonal
matrices in any chosen basis.  Its coordinate linear functionals,
after omitting zeros, split into proportionality classes.
A generic point of $S$ has even rank.  If $\dim S\ge2$, a generic
point on the hyperplane of one proportionality class and on none
of the others has rank lower by the size of that class.  That size
must therefore be even.  There are at most $n$ such classes,
and their representatives span $S^*$.  Thus $\dim S\le n$.
For $\dim S\le1$ the conclusion is immediate.

If $D_v$ is diagonal with distinct eigenvalues, then
\begin{equation}
 \kappa\lrcorner D_v
  =\sum_{j<k}(\lambda_{v,j}-\lambda_{v,k})
                      \kappa_{jk}q_{jk},
 \label{eq:tree-diagonal-tangent}
\end{equation}
where the $q_{jk}$ form the coordinate basis of $H^2(\mathcal O_A)$.
All displayed eigenvalue differences are nonzero.  Hence this map
on $\mathcal T$ has rank at least $n^2-n\ge2$.

\subsection{A leaf forces uniform multiplicities}

Choose a leaf orbit.  Its unique bridge is incident at one vertex
$b$ of that orbit.  For a fixed coordinate pair $j,k$, the internal
transposition edges form a matching, and their surface line is
$\mathbf Cq_{jk}$.  The coefficient in
\eqref{eq:tree-diagonal-tangent} has opposite signs at matched
vertices $v,(jk)v$.
If neither is $b$, their two trace equations therefore imply
\[
 (m_v-m_{(jk)v})(\lambda_{v,j}-\lambda_{v,k})\kappa_{jk}=0
                  \quad(\kappa\in\mathcal T).
\]
For every allowed coordinate pair, their multiplicities agree.

The allowed transpositions generate the symmetric group, since
the coordinate graph is connected.  Their Cayley graph on the
permutation orbit is connected and vertex transitive.  Deleting
one vertex leaves it connected: a leaf of a spanning tree is not
a cut vertex of a connected graph, and vertex transitivity then
shows that no vertex is a cut vertex.
Thus all vertices other than $b$ have a common multiplicity $m$.

Sum the trace equations over the leaf orbit.  Internal flows cancel.
The unweighted sum of its divisor matrices is zero, since every
eigenvalue occurs equally often in each position and their sum
is zero.  Thus the remaining bridge equation requires
\[
 (m_b-m)(\kappa\lrcorner D_b)\in\mathbf Cq_{\mathrm{bridge}}
                       \quad(\kappa\in\mathcal T).
\]
The image on the left before its scalar multiplication has
dimension at least two, whereas the right side has dimension one.
It follows that $m_b=m$, and the bridge flow is zero for every
$\kappa$.  The common scaling $N$ changes none of these conclusions.

Delete this leaf and its zero bridge and repeat.  Every orbit has
constant multiplicity.  On the final orbit the same conclusion
follows directly from the connected allowed-transposition graph.
The average \eqref{eq:finite-positive-average} now gives
\[
 \operatorname{ch}(\mathcal E)
  =\sum_a m_a(2n)!F_n(N\theta)\in\mathbf Q[\theta].
\]
A nonzero primitive middle Weil class is not a multiple of
$\theta^n$, so this contradicts $c\ne0$.
At least one Weil tangent therefore has nonzero Atiyah obstruction.
Its character obstruction is zero because the displayed character
stays Hodge.  The trace/evaluation identity sends that nonzero
obstruction to zero, proving failure of strong semiregularity.
This completes the proof of
Theorem~\ref{thm:tree-extension-obstruction}.

If the bundle has an additional coherent finite-group structure
whose action on degree-two Hochschild cohomology is trivial, the same
obstruction lies in the invariant Ext group, by equivariance of
evaluation.  It also excludes injectivity of the restricted
semiregularity map under that additional hypothesis.

The remaining construction question is more specific than simplicity:
one needs a gluing with enough independent bridges to satisfy all
the diagonal trace equations, and then the full second-Ext condition.
Solving the trace equations alone would be a necessary step and would
not establish semiregularity.
